\textbf{Phần 1. Trắc nghiệm lựa chọn}
\setcounterpageref{ex}{0}

\begin{ex} %[5X12N74]
Dữ liệu nào sau đây là dữ liệu định lượng (bằng số)?
\choice
{Xếp loại rèn luyện: Tốt, Khá, Đạt}
{\True Chiều cao các học sinh: $135\text{ cm}, 140\text{ cm}, 138\text{ cm}$}
{Môn học yêu thích: Toán, Tiếng Việt, Tiếng Anh}
{Màu sắc bóng: Xanh, Đỏ, Vàng}
\loigiai{
Chiều cao là dữ liệu định lượng bằng số.\\
\textbf{Đáp án của câu hỏi là B}
}
\end{ex}

\begin{ex} %[5X12N74]
Biểu đồ quạt tròn biểu diễn một hình quạt chiếm một góc vuông ($90^\circ$) tương ứng với bao nhiêu phần trăm?
\choice
{$20\%$}
{$30\%$}
{$50\%$}
{\True $25\%$}
\loigiai{
$\dfrac{90}{360} = \dfrac{1}{4} = 25\%$.\\
\textbf{Đáp án của câu hỏi là D}
}
\end{ex}

\begin{ex} %[5X12H74]
Gieo một con xúc xắc $60$ lần, mặt $3$ chấm xuất hiện $12$ lần. Tỉ số phần trăm số lần xuất hiện mặt $3$ chấm là
\choice
{\True $20\%$}
{$15\%$}
{$25\%$}
{$18\%$}
\loigiai{
$12 : 60 = 0{,}2 = 20\%$.\\
\textbf{Đáp án của câu hỏi là A}
}
\end{ex}

\begin{ex} %[5X12H74]
Biểu đồ quạt tròn biểu diễn thể loại sách mượn tại thư viện gồm: Truyện tranh $45\%$, Khoa học $30\%$, Lịch sử $25\%$. Biết có $200$ cuốn sách được mượn. Số cuốn sách khoa học được mượn là
\choice
{$50$ cuốn}
{$70$ cuốn}
{\True $60$ cuốn}
{$90$ cuốn}
\loigiai{
$200 \times 30 : 100 = 60$ (cuốn).\\
\textbf{Đáp án của câu hỏi là C}
}
\end{ex}

\begin{ex} %[5X12H74]
Cho bảng số liệu nhiệt độ $5$ ngày: $26^\circ\text{C}, 28^\circ\text{C}, 30^\circ\text{C}, 31^\circ\text{C}, 25^\circ\text{C}$. Nhiệt độ trung bình của $5$ ngày là
\choice
{\True $28^\circ\text{C}$}
{$27^\circ\text{C}$}
{$29^\circ\text{C}$}
{$30^\circ\text{C}$}
\loigiai{
Tổng nhiệt độ: $26 + 28 + 30 + 31 + 25 = 140^\circ\text{C}$.\\
Trung bình: $140 : 5 = 28^\circ\text{C}$.\\
\textbf{Đáp án của câu hỏi là A}
}
\end{ex}

\begin{ex} %[5X12H74]
Tung một đồng xu $80$ lần có $44$ lần xuất hiện mặt Ngửa. Tỉ số của số lần xuất hiện mặt Sấp so với tổng số lần tung là
\choice
{$\dfrac{11}{20}$}
{\True $\dfrac{9}{20}$}
{$\dfrac{9}{11}$}
{$\dfrac{11}{9}$}
\loigiai{
Mặt Sấp: $80 - 44 = 36$ lần.\\
Tỉ số: $\dfrac{36}{80} = \dfrac{9}{20}$.\\
\textbf{Đáp án của câu hỏi là B}
}
\end{ex}

\begin{ex} %[5X12V74]
Biểu đồ hình quạt tròn cho biết cơ cấu trái cây trong cửa hàng: Cam $40\%$, Xoài $35\%$, Táo $25\%$. Biết số táo là $50\text{ kg}$. Tổng khối lượng trái cây trong cửa hàng là
\choice
{$150\text{ kg}$}
{$250\text{ kg}$}
{$300\text{ kg}$}
{\True $200\text{ kg}$}
\loigiai{
Tổng khối lượng: $50 : 25 \times 100 = 200\text{ (kg)}$.\\
\textbf{Đáp án của câu hỏi là D}
}
\end{ex}

\begin{ex} %[5X12V74]
Một xạ thủ bắn súng $40$ lần, trong đó có $28$ lần trúng điểm 10, $8$ lần trúng điểm 9, còn lại trúng điểm 8. Tỉ số phần trăm số phát bắn không đạt điểm 10 là
\choice
{$25\%$}
{\True $30\%$}
{$35\%$}
{$70\%$}
\loigiai{
Số phát không đạt điểm 10: $40 - 28 = 12$ phát.\\
Tỉ số phần trăm: $12 : 40 = 30\%$.\\
\textbf{Đáp án của câu hỏi là B}
}
\end{ex}

\begin{ex} %[5X12V74]
Khảo sát số cây hoa trồng được của $4$ lớp: Lớp 5A ($45$ cây), Lớp 5B ($55$ cây), Lớp 5C ($50$ cây), Lớp 5D ($50$ cây). Tỉ số phần trăm số cây của lớp 5B so với tổng số cây của cả khối là
\choice
{\True $27{,}5\%$}
{$25\%$}
{$30\%$}
{$22{,}5\%$}
\loigiai{
Tổng số cây: $45 + 55 + 50 + 50 = 200$ cây.\\
Lớp 5B: $55 : 200 = 27{,}5\%$.\\
\textbf{Đáp án của câu hỏi là A}
}
\end{ex}

\begin{ex} %[5X12V74]
Biểu đồ hình quạt tròn thể hiện $4$ sở thích của học sinh: Bóng đá ($45\%$), Bơi lội ($25\%$), Cờ vua ($18\%$), Cầu lông ($12\%$). Số học sinh thích bóng đá nhiều hơn số học sinh thích bơi lội là $40$ em. Số học sinh tham gia khảo sát là
\choice
{$160$ em}
{$240$ em}
{\True $200$ em}
{$300$ em}
\loigiai{
Hiệu tỉ lệ: $45\% - 25\% = 20\%$.\\
Tổng số học sinh: $40 : 20 \times 100 = 200$ (em).\\
\textbf{Đáp án của câu hỏi là C}
}
\end{ex}

\begin{ex} %[5X12V74]
Trong một hộp có một số viên bi gồm $3$ màu: đỏ, xanh, vàng. Rút ngẫu nhiên $50$ lần thì thấy bi đỏ xuất hiện $22$ lần, bi xanh xuất hiện $18$ lần. Tỉ số phần trăm số lần rút được bi vàng là
\choice
{$15\%$}
{$25\%$}
{$30\%$}
{\True $20\%$}
\loigiai{
Bi vàng: $50 - (22 + 18) = 10$ lần.\\
Tỉ số phần trăm: $10 : 50 = 20\%$.\\
\textbf{Đáp án của câu hỏi là D}
}
\end{ex}

\begin{ex} %[5X12C74]
Điểm kiểm tra của $5$ học sinh là $6, 7, 8, 9, 10$. Nếu bổ sung thêm điểm kiểm tra của học sinh thứ sáu thì điểm trung bình của cả $6$ bạn tăng thêm $0{,}2$ điểm. Điểm của học sinh thứ sáu là
\choice
{$9$ điểm}
{$9{,}5$ điểm}
{\True $9{,}2$ điểm}
{$10$ điểm}
\loigiai{
Tổng điểm 5 bạn: $6 + 7 + 8 + 9 + 10 = 40$ điểm $\Rightarrow$ Điểm trung bình $= 8$ điểm.\\
Điểm trung bình 6 bạn: $8 + 0{,}2 = 8{,}2$ điểm.\\
Tổng điểm 6 bạn: $8{,}2 \times 6 = 49{,}2$ điểm.\\
Điểm bạn thứ 6: $49{,}2 - 40 = 9{,}2$ điểm.\\
\textbf{Đáp án của câu hỏi là C}
}
\end{ex}

\textbf{Phần 2. Trắc nghiệm đúng sai}
\setcounterpageref{ex}{0}

\begin{ex} %[5X12H74]
Bảng khảo sát món ăn sáng của $100$ học sinh: Bánh mì ($40$ bạn), Xôi ($30$ bạn), Phở ($20$ bạn), Bún ($10$ bạn). Xét tính đúng sai:
\choiceTF
{\True Học sinh thích Bánh mì chiếm $40\%$.}
{\True Học sinh thích Xôi chiếm $30\%$.}
{Số học sinh thích Phở và Bún chiếm $40\%$ tổng số học sinh.}
{\True Số học sinh thích Phở và Bún chiếm $30\%$ tổng số học sinh.}
\loigiai{
\begin{itemize}
    \item Bánh mì: $40 : 100 = 40\%$. Đúng.
    \item Xôi: $30 : 100 = 30\%$. Đúng.
    \item Phở + Bún: $20 + 10 = 30$ bạn ($30\% \neq 40\%$). Ý 3 Sai, Ý 4 Đúng.
\end{itemize}
\begin{itemize}
    \item a) \textbf{Đúng}.
    \item b) \textbf{Đúng}.
    \item c) \textbf{Sai}.
    \item d) \textbf{Đúng}.
\end{itemize}
}
\end{ex}

\begin{ex} %[5X12V74]
Gieo một con xúc xắc $100$ lần. Kết quả: Mặt chẵn ($54$ lần), mặt lẻ ($46$ lần). Trong đó mặt 6 chấm xuất hiện $18$ lần, mặt 1 chấm xuất hiện $14$ lần. Xét tính đúng sai:
\choiceTF
{\True Tỉ số phần trăm mặt chẵn xuất hiện là $54\%$.}
{\True Tỉ số phần trăm mặt 6 chấm xuất hiện là $18\%$.}
{Mặt lẻ xuất hiện nhiều hơn mặt chẵn.}
{\True Tỉ số phần trăm mặt lẻ xuất hiện là $46\%$.}
\loigiai{
\begin{itemize}
    \item Mặt chẵn: $54\%$. Đúng.
    \item Mặt 6 chấm: $18\%$. Đúng.
    \item Mặt lẻ ($46$ lần) ít hơn mặt chẵn ($54$ lần). Sai.
    \item Mặt lẻ: $46\%$. Đúng.
\end{itemize}
\begin{itemize}
    \item a) \textbf{Đúng}.
    \item b) \textbf{Đúng}.
    \item c) \textbf{Sai}.
    \item d) \textbf{Đúng}.
\end{itemize}
}
\end{ex}

\textbf{Phần 3. Trả lời ngắn}
\setcounterpageref{ex}{0}

\begin{ex} %[5X12H74]
Tung đồng xu $50$ lần thấy xuất hiện mặt Sấp $23$ lần. Tính số lần xuất hiện mặt Ngửa.
\shortans{27}
\loigiai{
\textbf{Đáp số:} $27$\\
$50 - 23 = 27$.
}
\end{ex}

\begin{ex} %[5X12H74]
Biểu đồ quạt tròn cho biết học sinh giỏi chiếm $35\%$, học sinh khá chiếm $45\%$, còn lại là học sinh trung bình. Tính tỉ số phần trăm học sinh trung bình. (Chỉ ghi số).
\shortans{20}
\loigiai{
\textbf{Đáp số:} $20$\\
$100 - (35 + 45) = 20\%$.
}
\end{ex}

\begin{ex} %[5X12V74]
Một cửa hàng bán được $400\text{ kg}$ gạo gồm $3$ loại: Gạo Bắc Hương ($40\%$), Gạo ST25 ($35\%$), còn lại là Gạo Tạp Giao. Tính khối lượng gạo Tạp Giao đã bán theo ki-lô-gam.
\shortans{100}
\loigiai{
\textbf{Đáp số:} $100$\\
Tỉ lệ Tạp Giao: $100\% - (40\% + 35\%) = 25\%$.\\
$400 \times 25 : 100 = 100\text{ kg}$.
}
\end{ex}

\begin{ex} %[5X12V74]
Điểm kiểm tra môn Toán của $4$ bạn: $8, 9, 9, 10$. Tính điểm trung bình của $4$ bạn.
\shortans{9}
\loigiai{
\textbf{Đáp số:} $9$\\
$(8 + 9 + 9 + 10) : 4 = 36 : 4 = 9$.
}
\end{ex}

\begin{ex} %[5X12V74]
Một nhóm sản xuất kiểm tra $250$ bóng đèn thấy có $10$ bóng bị hỏng. Tính tỉ số phần trăm số bóng đèn đạt chất lượng. (Chỉ ghi số).
\shortans{96}
\loigiai{
\textbf{Đáp số:} $96$\\
Bóng tốt: $240 \Rightarrow 240 : 250 = 96\%$.
}
\end{ex}

\begin{ex} %[5X12C74]
Biểu đồ quạt tròn cho biết tỉ lệ bình chọn tiết mục xuất sắc: Đội A ($35\%$), Đội B ($40\%$), Đội C ($25\%$). Biết Đội B nhận được nhiều hơn Đội C là $30$ phiếu bầu. Tính tổng số phiếu bầu của cả 3 đội.
\shortans{200}
\loigiai{
\textbf{Đáp số:} $200$\\
Hiệu tỉ lệ: $40\% - 25\% = 15\%$.\\
Tổng số phiếu: $30 : 15 \times 100 = 200$ phiếu.
}
\end{ex}
